\begin{abstract}
The \textsc{Apeiron} Framework proposes a non-anthropocentric approach to knowledge genesis, grounded in a foundational layer of irreducible observables defined via multi-axial atomicity. This paper presents an empirical implementation of that layer and evaluates its capacity for unsupervised structure discovery. We introduce a working Python implementation of the \texttt{UltimateObservable} class and its four decomposition operators (Boolean, measure-theoretic, categorical, and information-theoretic). We then apply the framework to the MNIST dataset, constructing a weighted hypergraph from raw pixel co-activations and identifying functional clusters via Louvain community detection. The resulting clusters exhibit significant modularity ($Q = 0.42$) and their ablation substantially degrades downstream classification accuracy ($p < 0.001$). We compare our method against standard baselines ($k$-means and spectral clustering) and find that the Apeiron-derived clusters align more closely with task-relevant spatial regions. While the atomicity concept itself is not fully exercised in this initial probe, the results demonstrate that the framework can autonomously extract functionally relevant structure from unlabeled data. We discuss the limitations of the current study and outline a roadmap for scaling the validation to more complex domains.
\end{abstract}